\paragraph{Interaction model.}
The user-facing interface supports conversational sessions with persistent thread history.
Within a session, requests may alternate between high-level scientific questions (``what data are available for mixed-layer depth?'') and procedural instructions (``plot a time series in this region and period'').
Responses may include explanatory text, generated figures, downloadable data subsets, and optional executable code.

\paragraph{Artifact presentation.}
Artifacts are treated as first-class outputs.
Figures can be previewed inline, while data subsets are offered as downloadable files suitable for direct use in notebooks or scripts.
This presentation helps bridge exploratory dialogue with reproducible analysis workflows.

\paragraph{Traceability.}
When tool-based operations are executed, provenance information (e.g., selected datasets/variables, bounds, and processing steps) can be summarized alongside results.
In practice, the interface can present this provenance as an optional, compact view (e.g., an expandable tool trace that can be copied or archived).
Such traces are important for scientific credibility and for debugging multi-step requests, consistent with broader calls for grounded, interpretable agent trajectories \citep{Yao2022ReAct}.
